\documentclass[11pt,a4paper]{article}
\usepackage[margin=2.4cm]{geometry}
\usepackage[T1]{fontenc}
\usepackage{lmodern,microtype}
\usepackage{amsmath,amssymb,amsthm,booktabs,array,enumitem,xcolor}
\usepackage{algorithm,algpseudocode}
\usepackage[colorlinks,allcolors=blue!55!black]{hyperref}
\newtheorem{theorem}{Theorem}
\newtheorem{proposition}{Proposition}
\theoremstyle{definition}
\newtheorem{definition}{Definition}
\theoremstyle{remark}
\newtheorem{remark}{Remark}
\newcommand{\nul}{\varnothing}
\newcommand{\cS}{\mathcal{S}}
\newcommand{\cG}{\mathcal{G}}
\newcommand{\Top}{\mathrm{Top}}
\DeclareMathOperator*{\argmax}{arg\,max}
\DeclareMathOperator*{\argmin}{arg\,min}
\setlist{itemsep=2pt,topsep=4pt}
\newcolumntype{L}[1]{>{\raggedright\arraybackslash}p{#1}}
\title{Graded Grounding with Priced Tree Answers\\[4pt]
\large Problem, algorithm and mathematics --- working document}
\author{Extension of \emph{Ask, Don't Draw} (ADD)}
\date{Version of 23 November 2026}

\begin{document}
\maketitle

\noindent\fbox{\parbox{0.97\linewidth}{\textbf{Status of this document.} This is a design document. No experiment
has been run for the method described here. Every proof below is complete and written out; the recursions,
Theorem~\ref{thm:nested} and Propositions~\ref{prop:logz}--\ref{prop:limit} were also checked numerically against
brute-force enumeration (Sec.~\ref{sec:status}). Numbers in the examples are illustrative unless they are marked as
results of the ADD paper. Statements about other papers are marked in the reference list as either read for this
document or recalled from memory.}}

\tableofcontents

%=====================================================================================================
\section{Summary}
\paragraph{Motivation.}
Every grounding model today is forced to commit. Given a query, it returns one region or it returns nothing, and
both answers look equally confident. But references in the real world are graded. ``The handle of the mug'' may
point to one handle, to three, or to a mug that has no visible handle; the model may be sure about the object and
unsure about the part. A model that can only commit or refuse must guess in all of these cases, and its guesses
fail silently. Recent probes show how common this is: when an expected part is removed from an image, none of ten
current vision--language models recognises its absence in more than 75\% of cases \cite{missing}, and on queries
that should be rejected most models score 0\% \cite{groundingme}.

We argue that grounding should be \emph{graded}: a model should answer at the finest level of detail it can stand
behind, and that level should come with a guarantee. The part when it is sure of the part; the object when the part
is uncertain or not visible; several regions when the reference is ambiguous; nothing when the object is absent.
This turns grounding from a single-answer prediction into a decision about specificity, in which a coarser answer is
a valid outcome and no longer counts as an error.

The obstacle is that such answers live on a hierarchy that is different for every image and query. Existing
methods with guarantees on hierarchies fix the hierarchy in advance, restrict answers to a single chain, or lose
the monotonicity that guarantees require (Sec.~\ref{sec:prior}). We show that one formulation removes this
obstacle: price the area an answer claims. The optimal answers are then nested as the price changes, which makes
the guarantee available, and the same rule yields a training loss.

\paragraph{Approach.}
We build on ADD \cite{add}, which scores about 260 candidate masks per query and arranges them in a containment
tree. We keep its candidates, tree and scores, and replace only the step that turns scores into an answer. The new
answer is a small \emph{set of nodes of the containment tree}. It is chosen by one rule: take the regions whose
probability is worth more than their area, at a price $\tau$ per unit of area. We show three things.
\begin{enumerate}
\item The answers for different prices are nested (Theorem~\ref{thm:nested}). A lower price can only enlarge the
answer.
\item If the probabilities are correct, the rule is optimal: no other rule covers as many targets with less claimed
area (Theorem~\ref{thm:optimal}).
\item Because of nesting, the price can be calibrated on held-out data so that the fraction of missed targets is at
most a chosen level $\alpha$, for any model (Theorem~\ref{thm:crc}, which is the conformal risk control theorem of
\cite{crc}).
\end{enumerate}
We then derive a training loss from the same rule (Sec.~\ref{sec:training}). It is the negative log-probability that a
smoothed version of the rule returns an acceptable answer, and it becomes the exact rule at zero temperature.
Finally we define a new task, \emph{graded grounding}, in which the correct answer can be a part, several parts, the
object when the part is not visible, or nothing (Sec.~\ref{sec:task}).

%=====================================================================================================
\section{The problem}
\subsection{What ADD does and where it stops}
For an image $I$ and a query $x$, ADD (our CVPR 2027 submission \cite{add}) performs four steps.
\begin{enumerate}[label=\textbf{\Alph*.}]
\item \textbf{Candidates.} A frozen segmenter (SAM\,3 \cite{sam3}) proposes masks from a point grid and from the
phrases of the query. This gives candidates $V=\{1,\dots,N\}$ with masks $m_v\subseteq\Omega$, where $\Omega$ is the
set of pixels and $N\approx260$.
\item \textbf{Tree.} The parent of a mask is the smallest larger mask that contains at least a fraction
$\rho=0.8$ of it. The root $r$ is the whole image, $m_r=\Omega$.
\item \textbf{Scores.} Each mask becomes one token. One forward pass of the multimodal language model gives a score
\[
s_v=\lambda\cos\!\big(\psi_q(h_{\langle\mathrm{SEL}\rangle}),\psi_k(z_v)\big)+\beta .
\]
\item \textbf{Answer.} $v^\star=\argmax_v s_v$; the answer is $m_{v^\star}$, or nothing if $\sigma(s_{v^\star})<\tau_{\mathrm{rej}}$.
\end{enumerate}
Steps A--C stay unchanged in everything that follows. We replace step~D.

Step~D has three weaknesses.
\begin{itemize}
\item \textbf{No signal of uncertainty.} A wrong mask looks like a right one. On PARTONOMY, ADD reaches 49.4 gIoU
while the best candidate in its pool reaches 71.0 \cite{add}; many answers are wrong although a good candidate
exists.
\item \textbf{Only two kinds of answer.} One mask, or nothing. There is no way to answer ``one of these two
handles'' or ``somewhere on this mug''.
\item \textbf{No guarantee.} The rejection threshold is tuned for a benchmark score. It says nothing about how often
the answer misses the target.
\end{itemize}

\subsection{The new task: graded grounding}
\label{sec:task}
Existing benchmarks have two kinds of query: the target exists, or it does not (gRefCOCO \cite{gres}). We consider
five cases. Example: an image with three mugs.

\begin{center}
\begin{tabular}{L{3.3cm}L{5.0cm}L{6.0cm}}
\toprule
Case & Example query & Correct answer \\
\midrule
Unique part & the handle of the red mug & that handle \\
All instances & the handles of the mugs & all three handles \\
Ambiguous & the handle of the mug & all three, or a hedge that contains them \\
Part not visible & the lid of the mug (no lid visible) & the mug, with the statement ``no lid visible'' \\
Object absent & the handle of the teapot & nothing \\
\bottomrule
\end{tabular}
\end{center}

\noindent The fourth case is the new one: the correct answer lies at an intermediate level of the hierarchy. The
goal is to return, for every query, \emph{the finest answer that is correct}, together with a guarantee on how often
the answer misses a target.

\paragraph{How the task can be built.} PACO-LVIS \cite{paco} annotates, for each object instance, the masks of its
visible parts from a fixed list of possible parts per object class, and it flags the images in which an object-part
is not exhaustively annotated. On the exhaustively annotated subset, a part of the class list without a mask on an
instance gives a ``part not visible'' query, and images with several instances of one class give the ``all
instances'' and ``ambiguous'' queries. No new labelling is needed. (``Not visible'' includes occluded parts; for
grounding, both have the same correct answer.)

%=====================================================================================================
\section{Notation}
\begin{center}
\begin{tabular}{ll}
\toprule
Symbol & Meaning \\
\midrule
$\Omega$ & pixels of the image \\
$V$, $m_v$ & candidates and their masks; $r\in V$ is the image root \\
$D(v)$ & $v$ together with all its descendants in the tree \\
$R(v)$ & region of $v$ (Def.~\ref{def:region}) \\
$\nul$ & extra leaf meaning ``no target'' \\
$s_v$, $q_v$ & score and probability of candidate $v$ (Sec.~\ref{sec:prob}) \\
$\gamma_v$ & own area of $v$ (Def.~\ref{def:cost}); $\gamma_\nul=c_\nul$ \\
$S$ & an answer: a downward-closed set of nodes (Def.~\ref{def:answer}) \\
$P(S)$, $C(S)$ & mass and cost of an answer \\
$\tau$ & price per unit of area \\
$w_u=q_u-\tau\gamma_u$ & weight of node $u$; $W(v)=\sum_{u\in D(v)}w_u$ \\
$\kappa$ & fraction of a target that must be covered \\
$\alpha$ & tolerated miss rate \\
$T$ & temperature of the smoothed rule \\
\bottomrule
\end{tabular}
\end{center}

%=====================================================================================================
\section{Inference: from scores to an answer}
\subsection{Component 1: regions}
In ADD's tree a child lies only 80\% inside its parent. Backing off from a child to its parent could therefore lose
pixels of the target. We remove this by giving each node the union of all masks below it.

\begin{definition}[Region]\label{def:region}
$R(v)=\bigcup_{u\in D(v)}m_u$ for $v\in V$.
\end{definition}

\begin{proposition}[Exact nesting]\label{prop:nest}
If $u\in D(v)$ then $R(u)\subseteq R(v)$.
\end{proposition}
\begin{proof}
$u\in D(v)$ implies $D(u)\subseteq D(v)$, so every mask in the union defining $R(u)$ also appears in the union
defining $R(v)$.
\end{proof}
\noindent Consequence: if a target is covered by $R(u)$, it is covered by the region of every ancestor of $u$.
Moving up the tree can make an answer coarser but never wrong.

\subsection{Component 2: probabilities}
\label{sec:prob}
\begin{definition}[Probabilities]
For every candidate $v\in V$, $q_v=\sigma(s_v)=1/(1+e^{-s_v})$.
\end{definition}
$q_v$ is the probability that candidate $v$ is a target. The probabilities do not sum to one over the pool: each
answers its own yes-or-no question. This is how ADD is trained (a focal loss \cite{focal} on each $\sigma(s_v)$
separately), and it is what allows several targets. Their sum estimates the number of targets.

\paragraph{No target.} The sigmoids have no left-over mass for ``nothing there'', so its probability is defined
separately. With a frozen ADD model we use
\[
q_\nul=1-\max_{v\in V}q_v ,
\]
which reproduces ADD's rejection rule exactly ($\sigma(s_{v^\star})<\tau_{\mathrm{rej}}\iff q_\nul>1-\tau_{\mathrm{rej}}$);
this rule gives 94.0 / 90.9 N-acc on gRefCOCO \cite{add}. When the model is trained with the loss of
Sec.~\ref{sec:training}, we use a learned value $q_\nul=\sigma(s_\nul)$ with $s_\nul$ a linear read-out of
$h_{\langle\mathrm{SEL}\rangle}$.

\paragraph{Extended tree.} We add $\nul$ as a leaf of its own, next to the image root. The node set is
$V^+=V\cup\{\nul\}$; it is a forest with the two roots $r$ and $\nul$.

\begin{remark}
The guarantee of Sec.~\ref{sec:calib} holds for any values of $q$. Accurate probabilities make the answers more
specific, not more valid. If needed, the scores can be rescaled on validation data, $q_v=\sigma(a\,s_v+b)$.
\end{remark}

\subsection{Component 3: answers and their cost}
\begin{definition}[Answer]\label{def:answer}
An answer is a set $S\subseteq V^+$ that is closed downwards: $v\in S\Rightarrow D(v)\subseteq S$. The set of all
answers is $\cS$. $\Top(S)$ is the set of nodes of $S$ whose parent is not in $S$. The claimed region is
$A(S)=\bigcup_{v\in\Top(S)\cap V}R(v)$, and the answer asserts ``no target'' as a possibility if $\nul\in S$.
\end{definition}
Examples: $\{$handle 1$\}$; $\{$handle 1, handle 2$\}$; $D(\text{mug})$ shown as $\{$mug$\}$; $\{\nul\}$; $V^+$
(``cannot tell''). ADD's answers are the special cases $D(v^\star)$ and $\{\nul\}$.

\begin{proposition}[Lattice]\label{prop:lattice}
If $S,S'\in\cS$ then $S\cup S'\in\cS$ and $S\cap S'\in\cS$.
\end{proposition}
\begin{proof}
Let $v\in S\cup S'$, say $v\in S$. Then $D(v)\subseteq S\subseteq S\cup S'$. Let $v\in S\cap S'$. Then
$D(v)\subseteq S$ and $D(v)\subseteq S'$.
\end{proof}

\begin{definition}[Mass and cost]\label{def:cost}
Each pixel is owned by the smallest candidate mask that covers it (ties broken by a fixed order). The own area of
$u\in V$ is $\gamma_u=\#\{\text{pixels owned by }u\}/|\Omega|$, and $\gamma_\nul=c_\nul>0$ is a constant. For an
answer $S$,
\[
P(S)=\sum_{u\in S}q_u,\qquad C(S)=\sum_{u\in S}\gamma_u .
\]
\end{definition}
$P(S)$ is the expected number of targets among the nodes of $S$. $C(S)$ is the fraction of the image owned by the
nodes of $S$; it is at most the area of $A(S)$ and equal to it when sibling masks do not overlap. Both are sums
over nodes with non-negative terms; this additivity is what the theory uses. The number of nodes of $\Top(S)$ would
not be additive, which is why we do not use it as a cost. $c_\nul$ is a convention (for example the typical area of
a target); without it, adding $\nul$ to every answer would be free.

\subsection{Component 4: the rule}
\begin{definition}[Rule]\label{def:rule}
For a price $\tau\ge0$, let $w_u=q_u-\tau\gamma_u$ and $w(S)=\sum_{u\in S}w_u=P(S)-\tau C(S)$. The answer is
\[
S(\tau)=\text{the largest maximiser of } w(S)\text{ over }S\in\cS .
\]
\end{definition}
The largest maximiser exists and is unique: by Proposition~\ref{prop:lattice} and additivity of $w$, the union of
two maximisers is again a maximiser (this is the case $\tau=\tau'$ of the proof of Theorem~\ref{thm:nested}).

\paragraph{Computation.} For a node $v$, an answer restricted to the subtree of $v$ either contains $v$, and then
it is all of $D(v)$, or it does not contain $v$, and then it is a union of independent answers in the subtrees of
the children. Hence, with $g(v)$ the best value inside the subtree of $v$,
\begin{equation}\label{eq:dp}
g(v)=\max\Big(W(v),\ \sum_{c\in\mathrm{ch}(v)}g(c)\Big),\qquad g(\text{leaf})=\max(w_{\text{leaf}},0).
\end{equation}
If the first term is at least as large as the second, $v$ is taken with its whole subtree and becomes a node of
$\Top(S)$; otherwise $v$ is left out and the children decide. Taking the subtree on ties returns the largest
maximiser. The total is $g(r)+g(\nul)$. Each node is visited once, so the cost is $O(N)$ (Algorithm~\ref{alg:rule}).

\paragraph{When a parent replaces its children.} $W(v)\ge\sum_cg(c)$ is equivalent to
$w_v\ge\sum_c\big(g(c)-W(c)\big)$. The right-hand side is the value of what the children had left out. A parent takes
over when its own probability pays for its own area and for everything its children had rejected.

\begin{theorem}[Nested answers]\label{thm:nested}
If $\tau<\tau'$ then $S(\tau')\subseteq S(\tau)$.
\end{theorem}
\begin{proof}
Write $w_\tau(S)=P(S)-\tau C(S)$. Let $S$ be a maximiser at $\tau$ and $S'$ a maximiser at $\tau'$, and let
$\Delta=S'\setminus S$. By Proposition~\ref{prop:lattice}, $S\cap S'\in\cS$ and $S\cup S'\in\cS$. Since $S'$ is optimal
at $\tau'$ and $S'=(S\cap S')\cup\Delta$ with a disjoint union,
\[
0\le w_{\tau'}(S')-w_{\tau'}(S\cap S')=P(\Delta)-\tau'C(\Delta).
\]
Because $C(\Delta)\ge0$ and $\tau<\tau'$, also $P(\Delta)-\tau C(\Delta)\ge0$. Therefore
$w_\tau(S\cup S')=w_\tau(S)+P(\Delta)-\tau C(\Delta)\ge w_\tau(S)$, so $S\cup S'$ is a maximiser at $\tau$. Taking $S$
to be the largest maximiser at $\tau$ gives $S\cup S'\subseteq S$, that is, $S'\subseteq S(\tau)$ for every
maximiser $S'$ at $\tau'$, in particular for $S(\tau')$.
\end{proof}

\begin{theorem}[Optimality]\label{thm:optimal}
Suppose that for every input the values $q_u$ equal the true probabilities that node $u$ is a truth node
(Def.~\ref{def:truth}). Fix $\tau>0$ and let $S_\tau(\cdot)$ be the rule of Def.~\ref{def:rule}. Let $S'(\cdot)$ be
any rule that returns an answer in $\cS$ for every input and covers at least as many truth nodes in expectation,
$\mathbb{E}[P(S')]\ge\mathbb{E}[P(S_\tau)]$. Then $\mathbb{E}[C(S_\tau)]\le\mathbb{E}[C(S')]$.
\end{theorem}
\begin{proof}
For every input, $S_\tau$ maximises $P(S)-\tau C(S)$, so $P(S_\tau)-\tau C(S_\tau)\ge P(S')-\tau C(S')$. Taking
expectations and using $\mathbb{E}[P(S')]\ge\mathbb{E}[P(S_\tau)]$ gives $\tau\,\mathbb{E}[C(S_\tau)]\le\tau\,
\mathbb{E}[C(S')]$.
\end{proof}
\noindent The theorem is the tree version of the Neyman--Pearson argument \cite{np}: threshold probability per unit
of cost. For flat label sets the analogous result is known \cite{sadinle}.

\subsection{Component 5: calibration of the price}
\label{sec:calib}
\begin{definition}[Miss loss]\label{def:loss}
Let a query have targets with masks $m_1,\dots,m_k$ ($k\ge0$) and let $S$ be an answer. Target $t$ is covered if
$|m_t\cap A(S)|\ge\kappa|m_t|$.
\[
L(S)=\begin{cases}\dfrac{\#\{t:\ t\text{ not covered}\}}{k}&k\ge1,\\[8pt] \mathbf{1}[\nul\notin S]&k=0.\end{cases}
\]
\end{definition}
$L\in[0,1]$. If $S\subseteq S'$ then $A(S)\subseteq A(S')$ and $\nul\in S\Rightarrow\nul\in S'$, so
$L(S')\le L(S)$. Together with Theorem~\ref{thm:nested} and Proposition~\ref{prop:nest}:
\begin{equation}\label{eq:mono}
\tau<\tau'\ \Longrightarrow\ L\big(S(\tau)\big)\le L\big(S(\tau')\big)\quad\text{for every single query,}
\end{equation}
and $L(S(0))=0$ because $S(0)=V^+$.

\paragraph{Procedure.} Take $n$ held-out images with ground truth. Let $L_i(\tau)$ be the loss of image $i$ at price
$\tau$, averaged over the queries of that image, and $\hat R(\tau)=\frac1n\sum_iL_i(\tau)$. On a fixed finite grid
$\Gamma\ni0$ of prices,
\begin{equation}\label{eq:tauhat}
\hat\tau=\max\Big\{\tau\in\Gamma:\ \frac{n}{n+1}\hat R(\tau)+\frac{1}{n+1}\le\alpha\Big\}.
\end{equation}
The set is not empty if $\alpha\ge1/(n+1)$, because $\hat R(0)=0$. Images, not queries, are the units, because
several queries on one image are dependent.

\begin{theorem}[Risk control {\cite{crc}}]\label{thm:crc}
If the $n$ calibration images and the test image are exchangeable, then
$\ \mathbb{E}\big[L_{n+1}(\hat\tau)\big]\le\alpha .$
\end{theorem}
\noindent This is the conformal risk control theorem of Angelopoulos et al.\ \cite{crc}, applied with the parameter
$-\tau$. Its conditions are: a loss bounded by $1$, monotone in the parameter (Eq.~\ref{eq:mono}), and a parameter
value at which the loss is at most $\alpha$ ($\tau=0$). On a finite grid no continuity condition is needed. We state
the theorem from the literature and do not re-prove it; its exact conditions should be checked against \cite{crc}
before publication. The expectation is over the calibration set and the test image: the guarantee is an average,
not a statement about each single answer.

\paragraph{The floor.} Some targets have no usable candidate. If the only node that covers such a target is the image
root, and the root has probability near zero, the rule takes it only at a price near zero. The miss rate at any
useful price is therefore at least the fraction of such targets. If $\alpha$ is below this floor, $\hat\tau$ is
close to $0$ and every answer is ``everything'': valid but empty. The floor is measurable without a model and is the
first thing to measure.

%=====================================================================================================
\section{Training through the rule}
\label{sec:training}
A frozen ADD model has a specific weakness for this task. For the query ``the handle of the mug'', the mug itself
was always a negative. The model has no way to put probability on the mug when no handle is visible or when the
pool contains no handle candidate. The loss below rewards exactly that.

\subsection{The gold answer}
\begin{definition}[Truth node, gold answer, ignored nodes]\label{def:truth}
For a target with mask $m_t$, the truth node is the smallest node that covers it,
\[
u_t=\argmin_{v\in V:\ |m_t\cap R(v)|\ge\kappa|m_t|}\ |R(v)| ,
\]
which exists because the root covers everything. The gold answer is $S^\star=\bigcup_tD(u_t)$ for a query with
targets and $S^\star=\{\nul\}$ for a query without. The ignored set $I\subseteq V\setminus S^\star$ contains the
candidates that match some target with IoU $\ge\theta$ but are not in $S^\star$ (for example a near-duplicate mask),
as in ADD. The acceptable answers are
\[
\cG=\{S\in\cS:\ S^\star\subseteq S\subseteq S^\star\cup I\}.
\]
\end{definition}
\begin{center}
\begin{tabular}{L{5.2cm}L{4.6cm}L{4.4cm}}
\toprule
Case & Truth node & Gold answer \\
\midrule
Part visible, good candidate & that candidate & the part \\
Part visible, no good candidate & smallest covering node, often the object & the object \\
Part not visible & smallest node covering the object & the object \\
Several targets & one per target & all of them \\
No target & --- & $\{\nul\}$ \\
\bottomrule
\end{tabular}
\end{center}
For ``part not visible'' queries the target mask is the mask of the object. Containing the truth node implies that
the target is covered in the sense of Def.~\ref{def:loss}; the converse can fail when several nodes together cover a
target that none covers alone.

\subsection{The smoothed rule}
Fix a training price $\tau_0$, weights $w_u=q_u-\tau_0\gamma_u$, and a temperature $T>0$. Define a distribution over
answers,
\begin{equation}\label{eq:gibbs}
p(S)=\frac{\exp(w(S)/T)}{Z},\qquad Z=\sum_{S\in\cS}\exp(w(S)/T).
\end{equation}
As $T\to0$ all probability goes to the maximisers of $w$, that is, to the rule of Def.~\ref{def:rule}.

\begin{proposition}[Partition function]\label{prop:logz}
Let $A(v)$ be defined bottom-up by
\[
A(\text{leaf})=\log\big(1+e^{w_{\text{leaf}}/T}\big),\qquad
A(v)=\log\Big(e^{W(v)/T}+\exp\!\Big(\sum_{c\in\mathrm{ch}(v)}A(c)\Big)\Big).
\]
Then $\log Z=A(r)+A(\nul)$.
\end{proposition}
\begin{proof}
Let $Z(v)$ be the sum of $\exp(w(S)/T)$ over the answers $S$ inside the subtree of $v$. Such an answer either
contains $v$, and then equals $D(v)$ with weight $W(v)$, or does not contain $v$, and then is a union of one answer
per child subtree, chosen independently; the weights add, so the exponentials multiply. Hence
$Z(v)=e^{W(v)/T}+\prod_cZ(c)$, with the empty product equal to $1$ for a leaf. $A(v)=\log Z(v)$. The two trees with
roots $r$ and $\nul$ are independent, so $Z=Z(r)Z(\nul)$.
\end{proof}

\begin{proposition}[Inclusion probabilities]\label{prop:marg}
Let $\pi(v)=\exp\big(W(v)/T-A(v)\big)$. Then for every node $u$,
\[
m(u):=p(u\in S)=1-\prod_{v\in\mathrm{path}(u)}\big(1-\pi(v)\big)=T\,\frac{\partial\log Z}{\partial w_u},
\]
where $\mathrm{path}(u)$ is the set of nodes from the root of $u$'s tree down to $u$, both included.
\end{proposition}
\begin{proof}
Condition on the event that no strict ancestor of $v$ is in $S$. By the product structure in the proof of
Proposition~\ref{prop:logz}, the part of $S$ inside the subtree of $v$ is then distributed as in
Eq.~\eqref{eq:gibbs} restricted to that subtree, so $v\in S$ has probability $e^{W(v)/T}/Z(v)=\pi(v)$. Node $u$ is
\emph{not} in $S$ exactly if none of the nodes on its path is taken, which has probability
$\prod_{v\in\mathrm{path}(u)}(1-\pi(v))$. The second equality is the standard identity
$\partial\log Z/\partial w_u=\frac1T\sum_{S\ni u}p(S)$.
\end{proof}

\subsection{The loss}
\begin{proposition}[Constrained partition function]\label{prop:cons}
Let $B=V^+\setminus(S^\star\cup I)$ be the forbidden nodes. Define $A_\cG(v)$ by the recursion of
Proposition~\ref{prop:logz} with two restrictions: the term $e^{W(v)/T}$ is present only if $D(v)\cap B=\emptyset$,
and the term $\exp(\sum_cA_\cG(c))$ (which is $1$ for a leaf) is present only if $v\notin S^\star$. Then
$A_\cG(r)+A_\cG(\nul)=\log Z_\cG$ with $Z_\cG=\sum_{S\in\cG}\exp(w(S)/T)$.
\end{proposition}
\begin{proof}
The first restriction removes exactly the answers that contain a forbidden node, since taking $v$ whole adds
$D(v)$. The second removes exactly the answers that omit a node of $S^\star$: a node $v\in S^\star$ is in $S$ iff it or
an ancestor is taken whole, and forbidding ``leave $v$ out'' at $v$ forces this whenever the recursion reaches $v$.
At least one term is always present: if $v\in S^\star$ then $D(v)\subseteq S^\star$, so $D(v)\cap B=\emptyset$.
\end{proof}

\begin{definition}[Rule loss]
\begin{equation}\label{eq:loss}
\mathcal{L}_{\mathrm{rule}}=-\log p(S\in\cG)=\log Z-\log Z_\cG .
\end{equation}
\end{definition}

\begin{proposition}[Gradient]\label{prop:grad}
With $m_\cG(u)=p(u\in S\mid S\in\cG)$,
\[
\frac{\partial\mathcal{L}_{\mathrm{rule}}}{\partial w_u}=\frac{m(u)-m_\cG(u)}{T},\qquad
\frac{\partial w_u}{\partial s_u}=q_u(1-q_u).
\]
\end{proposition}
\begin{proof}
Apply the identity of Proposition~\ref{prop:marg} to $\log Z$ and to $\log Z_\cG$.
\end{proof}
\noindent For a node of $S^\star$, $m_\cG=1$: the gradient raises its weight until the rule includes it. For a
forbidden node, $m_\cG=0$: the gradient lowers its weight until the rule excludes it. For an ignored node the two
terms largely cancel. The gradient reaches the mask tokens, the tree biases, the heads and the LoRA weights through
the scores.

\begin{proposition}[Zero-temperature limit]\label{prop:limit}
$\displaystyle\lim_{T\to0}T\,\mathcal{L}_{\mathrm{rule}}=\max_{S\in\cS}w(S)-\max_{S\in\cG}w(S)\ \ge0$, with equality iff
some maximiser of $w$ over $\cS$ is acceptable.
\end{proposition}
\begin{proof}
$T\log\sum_Se^{w(S)/T}\to\max_Sw(S)$ as $T\to0$ for a finite sum; apply this to $Z$ and $Z_\cG$, and note
$\cG\subseteq\cS$.
\end{proof}
\noindent So the loss measures by how much the test-time rule, at price $\tau_0$, prefers an unacceptable answer.
It has the form of the conditional likelihood of a conditional random field \cite{crf}, with the answers $\cS$ as
the structured output space.

\subsection{Full objective}
\[
\mathcal{L}=\mathcal{L}_{\mathrm{rule}}+\mathcal{L}_{\mathrm{ntp}}+\lambda\,\mathcal{L}_{\mathrm{sel}} .
\]
$\mathcal{L}_{\mathrm{ntp}}$ is ADD's next-token loss on the answer span, unchanged. $\mathcal{L}_{\mathrm{sel}}$ is
ADD's focal loss with the positives relabelled to the truth nodes; it is kept with a small weight because it is the
only term in which candidates of other images act as negatives. ADD's family loss is dropped, since
$\mathcal{L}_{\mathrm{rule}}$ already compares a node with its parent and children. Multi-target queries need no
special treatment: $S^\star$ is the union over the targets. ADD already trains on the multi-target expressions of
gRefCOCO with one positive per target.

%=====================================================================================================
\section{Algorithms}
\begin{algorithm}[H]
\caption{Answer for one query at price $\tau$}\label{alg:rule}
\begin{algorithmic}[1]
\Require image $I$, query $x$, price $\tau$, constants $c_\nul$
\State run ADD steps A--C: candidates $V$, tree, scores $s_v$
\State $R(v)\gets\bigcup_{u\in D(v)}m_u$ for all $v$ (one pass from the leaves)
\State $q_v\gets\sigma(s_v)$; \ $q_\nul\gets1-\max_vq_v$ (or $\sigma(s_\nul)$ for a trained model)
\State assign each pixel to the smallest mask covering it; $\gamma_v\gets$ owned fraction; $\gamma_\nul\gets c_\nul$
\State $w_u\gets q_u-\tau\gamma_u$ for all $u\in V^+$
\Function{Solve}{$v$}
  \State $(g_c,S_c,W_c)\gets$ \Call{Solve}{$c$} for each child $c$ of $v$
  \State $W\gets w_v+\sum_cW_c$; \ $g\gets\sum_cg_c$
  \If{$W\ge g$ \textbf{and} $W\ge0$} \Return $(W,\,D(v),\,W)$
  \Else\ \Return $(g,\,\bigcup_cS_c,\,W)$
  \EndIf
\EndFunction
\State $S\gets$ second output of \Call{Solve}{$r$} $\cup$ second output of \Call{Solve}{$\nul$}
\State \Return $\Top(S)$ and the region $A(S)$
\end{algorithmic}
\end{algorithm}

\begin{algorithm}[H]
\caption{Calibration of the price}\label{alg:calib}
\begin{algorithmic}[1]
\Require $n$ held-out images with ground truth, level $\alpha$, grid $\Gamma$ (about 100 prices, log-spaced, with $0$), $\kappa$
\For{each image $i$, each of its queries, each $\tau\in\Gamma$}
  \State compute $S(\tau)$ with Algorithm~\ref{alg:rule} and its loss (Def.~\ref{def:loss})
\EndFor
\State $L_i(\tau)\gets$ mean loss over the queries of image $i$; \ $\hat R(\tau)\gets\frac1n\sum_iL_i(\tau)$
\State \Return $\hat\tau\gets\max\{\tau\in\Gamma:\frac{n}{n+1}\hat R(\tau)+\frac1{n+1}\le\alpha\}$
\end{algorithmic}
\end{algorithm}

\begin{algorithm}[H]
\caption{One training step for one query}\label{alg:train}
\begin{algorithmic}[1]
\Require image, query, target masks (possibly none), $\tau_0$, $T$, $\kappa$, $\theta$, $\lambda$
\State run ADD steps A--C with gradients; compute $R$, $q$, $\gamma$, $w_u=q_u-\tau_0\gamma_u$
\State truth nodes $u_t$, gold answer $S^\star$, ignored set $I$ (Def.~\ref{def:truth}); $B\gets V^+\setminus(S^\star\cup I)$
\State $\log Z\gets A(r)+A(\nul)$ by the recursion of Proposition~\ref{prop:logz}
\State $\log Z_\cG\gets A_\cG(r)+A_\cG(\nul)$ by the recursion of Proposition~\ref{prop:cons}
\State $\mathcal{L}\gets(\log Z-\log Z_\cG)+\mathcal{L}_{\mathrm{ntp}}+\lambda\mathcal{L}_{\mathrm{sel}}$
\State back-propagate and update LoRA, the selection-token embedding, $f$, $\psi_q$, $\psi_k$, the tree biases and the $\nul$ read-out
\end{algorithmic}
\end{algorithm}
All recursions are evaluated in log space (log-sum-exp) for numerical stability and cost $O(N)$ per query.

%=====================================================================================================
\section{Worked example}
Query: ``the handles of all the mugs'', three mugs, the model unsure about mug~3. $c_\nul=0.01$. The numbers are
illustrative.
\begin{center}
\begin{tabular}{lcc}
\toprule
Node & $q$ & own area $\gamma$ \\
\midrule
handle 1 & 0.92 & 0.01 \\
handle 2 & 0.88 & 0.01 \\
handle 3 & 0.40 & 0.01 \\
rim of mug 3 & 0.35 & 0.01 \\
mug 3 (itself) & 0.10 & 0.06 \\
$\nul$ & $1-0.92=0.08$ & 0.01 \\
\bottomrule
\end{tabular}
\end{center}
At $\tau=10$: $w$(handle 3)$=0.30$, $w$(rim 3)$=0.25$, $w$(mug 3)$=0.10-0.60=-0.50$. The subtree of mug~3 has total
weight $0.05$, less than the $0.55$ of its two children, so the mug is not taken. $w(\nul)=-0.02<0$.
\begin{center}
\begin{tabular}{clcc}
\toprule
$\tau$ & Answer $\Top(S(\tau))$ & mass $P$ & cost $C$ \\
\midrule
50 & $\{$handle 1, handle 2$\}$ & 1.80 & 0.02 \\
10 & $\{$handle 1, handle 2, handle 3, rim 3$\}$ & 2.55 & 0.04 \\
1 & $\{$handle 1, handle 2, mug 3, $\nul\}$ & 2.73 & 0.10 \\
0 & everything & all & all \\
\bottomrule
\end{tabular}
\end{center}
The answers are nested, as Theorem~\ref{thm:nested} states (mug~3 contains handle~3 and rim~3). At $\tau=10$ the
answer for the third mug is ``handle 3 or rim 3'', a quarter of the area of the mug. Calibration picks the highest
price whose miss rate on held-out data is at most $\alpha$.

%=====================================================================================================
\section{Relation to prior work}
\label{sec:prior}
\begin{center}
\small
\begin{tabular}{L{3.0cm}L{5.5cm}L{6.3cm}}
\toprule
Our component & Prior work & Connection and difference \\
\midrule
Guarantee (Thm.~\ref{thm:crc}) & Conformal risk control \cite{crc} & We use their theorem unchanged. Their
hierarchical example backs off along the ancestors of the top prediction, a single chain. \\
Back-off on a hierarchy & Hedging Your Bets \cite{hedging} & Same idea of answering with the most specific node that
is accurate enough, on a fixed label taxonomy and without a finite-sample guarantee. Ours: a tree built per image
and query, spatial regions, sets of nodes, ``no target''. \\
Answers as sets of tree nodes & Conformal Structured Prediction \cite{csp}; representation complexity
\cite{mortier22,mortier26}; Hierarchical Conformal Classification \cite{hcc} & They choose at most $m$ nodes of a
tree or DAG under a hard constraint (required mass, or node cap). \cite{csp} notes that monotonicity in the
threshold is lost and uses learn-then-test with an integer program; \cite{hcc} states that no monotone loss is
known for its trade-off and uses a multiple-testing correction; \cite{mortier26} builds nested sets greedily. Ours
prices the cost instead of capping it; nesting then follows from the optimisation (Thm.~\ref{thm:nested}) and the
solver is linear-time. We do not cap the number of nodes. \\
Nesting imposed by hand & PAC prediction sets for code \cite{khakhar} & They enforce monotonicity as a constraint in
an integer program over syntax trees. Ours obtains it without a constraint. \\
Back-off with a guarantee & Conformal factuality \cite{mohri} & A single chain of less specific outputs. They also
show a coverage guarantee that does not need nesting. \\
Proof of Thm.~\ref{thm:nested} & Parametric closure and minimum cuts \cite{picard,ggt}; lattice programming
\cite{topkis} & The mathematical fact (nested solutions of a parametric closure problem) is classical. Our
contribution is the connection to prediction sets on trees. \\
Thm.~\ref{thm:optimal} & Neyman--Pearson \cite{np}; least ambiguous sets \cite{sadinle} & Same argument, for flat
sets. Ours: downward-closed sets of a per-input tree with an area cost. \\
Rule loss (Eq.~\ref{eq:loss}) & Conditional random fields \cite{crf} & Standard conditional likelihood; the
structured space (downward-closed sets) and its recursion are specific to our rule. \\
Per-candidate probabilities & Focal loss \cite{focal}; ADD \cite{add} & Unchanged from ADD. \\
Task: part not visible & MissingBench-Verified \cite{missing}; GroundingME \cite{groundingme}; gRefCOCO \cite{gres} &
\cite{missing} asks ``is the part visible?'' as multiple choice on 118 edited images, without grounding.
\cite{groundingme} has about 50 box-level queries with a wrong component where the correct answer is to reject.
\cite{gres} rejects whole expressions. None asks for the object as the answer when the part is not visible. \\
\bottomrule
\end{tabular}
\end{center}

\paragraph{What we would claim.} (i) The priced formulation, under which nesting and optimality hold together on
per-input trees; (ii) the loss derived from the rule; (iii) the graded grounding task. We would \emph{not} claim
hierarchical back-off, conformal guarantees on hierarchies, or input-dependent hierarchies as new.

\paragraph{Not yet checked.} Two lines of work have not been read for this document and must be before any claim of
novelty: set-based utility maximisation for set-valued prediction (cited in \cite{mortier22}), and referring
segmentation with false-premise queries.

%=====================================================================================================
\section{Status, parameters and risks}
\label{sec:status}
\paragraph{Proven in this document.} Propositions~\ref{prop:nest}--\ref{prop:limit}, Theorems~\ref{thm:nested} and
\ref{thm:optimal}.

\paragraph{Checked numerically.} \texttt{paper/icml/check\_math.py} compares, on 300 random trees with up to nine
nodes, against brute-force enumeration of all answers: the recursion~\eqref{eq:dp} returns the largest maximiser;
answers are nested over a grid of prices; $\log Z$ (error $<10^{-14}$); the inclusion probabilities of
Proposition~\ref{prop:marg} (error $<10^{-15}$) and the gradient identity (finite differences, error $<10^{-6}$);
the constrained $\log Z_\cG$; and the limit of Proposition~\ref{prop:limit}.

\paragraph{Taken from the literature.} Theorem~\ref{thm:crc}.

\paragraph{Not done.} No experiment. The coverage diagnostic, the construction of the graded queries, and any
training are still to do.

\paragraph{Parameters.}
\begin{center}
\begin{tabular}{lL{6.2cm}L{5.2cm}}
\toprule
Parameter & Role & How it is set \\
\midrule
$\kappa$ & fraction of a target that must be covered & from the coverage diagnostic on validation data \\
$c_\nul$ & price of answering ``nothing'' & convention, fixed on validation data \\
$\alpha$ & tolerated miss rate & chosen by the user; must exceed the floor \\
$\Gamma$, $n$ & price grid, number of calibration images & about 100 prices; a few hundred images suffice \\
$\tau_0$, $T$ & training price and temperature & to be tuned on validation data \\
$\lambda$ & weight of the auxiliary focal loss & small; to be tuned \\
$\theta$ & IoU for ignored nodes & $0.5$, as in ADD \\
\bottomrule
\end{tabular}
\end{center}

\paragraph{Risks.}
\begin{itemize}
\item \textbf{Calibration data.} The guarantee holds only for data exchangeable with the calibration set. Benchmarks
without a validation split have no guarantee unless part of the test set is used for calibration.
\item \textbf{The floor.} If many targets have no covering node other than the root, useful prices cannot reach low
miss rates.
\item \textbf{Frozen model.} Without training, hedging is mostly sideways (``handle 3 or rim 3''), because parents
carry little probability.
\item \textbf{The cost of $\nul$.} It is arbitrary and affects how often answers read ``\dots or nothing''.
\item \textbf{Training.} Possible collapse to large answers; sensitivity to $T$ and $\tau_0$; the simpler baseline
(relabelled positives with the focal loss only) may perform as well and must be compared.
\item \textbf{Novelty.} The mathematics is elementary and its core is classical; the contribution is the
formulation and what it removes. Two related lines are not yet checked (Sec.~\ref{sec:prior}).
\end{itemize}

%=====================================================================================================
\begin{thebibliography}{99}
\small
\item[] Marks: [R] read for this document; [M] recalled from memory, details to be verified.

\bibitem{add} [R] Anonymous. \emph{Ask, Don't Draw: Part-Level Grounding with Query-Conditioned Mask Tokens.}
Submission to CVPR 2027 (our own work).
\bibitem{crc} [R, theorem statement M] A.~N.~Angelopoulos, S.~Bates, A.~Fisch, L.~Lei, T.~Schuster.
\emph{Conformal Risk Control.} ICLR 2024. arXiv:2208.02814.
\bibitem{csp} [R] B.~Zhang, S.~Li, O.~Bastani. \emph{Conformal Structured Prediction.} ICLR 2025. arXiv:2410.06296.
\bibitem{mortier26} [R] T.~Mortier, A.~Javanmardi, Y.~Sale, E.~H\"ullermeier, W.~Waegeman. \emph{Conformal
Prediction in Hierarchical Classification with Constrained Representation Complexity.} AISTATS 2026.
arXiv:2501.19038.
\bibitem{mortier22} [R] T.~Mortier, E.~H\"ullermeier, K.~Dembczy\'nski, W.~Waegeman. \emph{Set-valued prediction in
hierarchical classification with constrained representation complexity.} UAI 2022. arXiv:2203.06676.
\bibitem{hcc} [R] F.~den~Hengst, I.~Blin, M.~Mohammadi, S.~I.~Shah, T.~Younesian. \emph{Hierarchical Conformal
Classification.} arXiv:2508.13288, 2025.
\bibitem{khakhar} [R] A.~Khakhar, S.~Mell, O.~Bastani. \emph{PAC Prediction Sets for Large Language Models of Code.}
ICML 2023. arXiv:2302.08703.
\bibitem{mohri} [R, venue M] C.~Mohri, T.~Hashimoto. \emph{Language Models with Conformal Factuality Guarantees.}
ICML 2024. arXiv:2402.10978.
\bibitem{hedging} [M] J.~Deng, J.~Krause, A.~C.~Berg, L.~Fei-Fei. \emph{Hedging Your Bets: Optimizing
Accuracy-Specificity Trade-offs in Large Scale Visual Recognition.} CVPR 2012.
\bibitem{sadinle} [M] M.~Sadinle, J.~Lei, L.~Wasserman. \emph{Least Ambiguous Set-Valued Classifiers with Bounded
Error Levels.} Journal of the American Statistical Association, 2019.
\bibitem{np} [M] J.~Neyman, E.~S.~Pearson. \emph{On the Problem of the Most Efficient Tests of Statistical
Hypotheses.} Philosophical Transactions of the Royal Society A, 1933.
\bibitem{picard} [M] J.-C.~Picard. \emph{Maximal Closure of a Graph and Applications to Combinatorial Problems.}
Management Science, 1976.
\bibitem{ggt} [M] G.~Gallo, M.~D.~Grigoriadis, R.~E.~Tarjan. \emph{A Fast Parametric Maximum Flow Algorithm and
Applications.} SIAM Journal on Computing, 1989.
\bibitem{topkis} [M] D.~M.~Topkis. \emph{Minimizing a Submodular Function on a Lattice.} Operations Research, 1978.
\bibitem{crf} [M] J.~Lafferty, A.~McCallum, F.~Pereira. \emph{Conditional Random Fields: Probabilistic Models for
Segmenting and Labeling Sequence Data.} ICML 2001.
\bibitem{focal} [M] T.-Y.~Lin, P.~Goyal, R.~Girshick, K.~He, P.~Doll\'ar. \emph{Focal Loss for Dense Object
Detection.} ICCV 2017.
\bibitem{paco} [R, venue M] V.~Ramanathan et al. \emph{PACO: Parts and Attributes of Common Objects.} CVPR 2023.
arXiv:2301.01795.
\bibitem{gres} [M] C.~Liu, H.~Ding, X.~Jiang. \emph{GRES: Generalized Referring Expression Segmentation.} CVPR 2023.
\bibitem{sam3} [M] N.~Carion et al. \emph{SAM 3: Segment Anything with Concepts.} 2026 (as cited in \cite{add}).
\bibitem{missing} [R] W.~M.~Guo, Q.~Qian, S.~Zhou, G.~Luo, S.~Du. \emph{MissingBench-Verified: Probing
Vision-Language Models' Inability to Detect Missing Object Parts.} arXiv:2607.18673, 2026.
\bibitem{groundingme} [R] R.~Li et al. \emph{GroundingME: Exposing the Visual Grounding Gap in MLLMs through
Multi-Dimensional Evaluation.} arXiv:2512.17495, 2025.
\end{thebibliography}
\end{document}
